% !TeX root = ../../main.tex
\chapter{补充理论与证明}[Supplementary Theory and Proofs]\label{cw:app:appendix}
\section{符号定义与说明}\label{cw:app:notation}

\textbf{概率单纯形与全一向量：}对正整数 $n$，$\mathbb R_+=[0,\infty)$，定义
\begin{equation*}
\Delta^{n-1}:=\left\{\bm w\in\mathbb R_+^n:\sum_{i=1}^n w_i=1\right\},\qquad
\mathbf1_n:=(1,\dots,1)^{\mathrm T}.
\end{equation*}
$\Delta^{n-1}$ 是 $n$ 个概率权重组成的集合，上标表示其维数；允许零权重，严格正权重需另行说明。

\textbf{Dirac 测度与示性函数：}对 $x\in\mathcal S$ 和可测集合 $A\subseteq\mathcal S$，
\begin{equation*}
\delta_x(A):=\mathbb I_A(x)=
\begin{cases}1,&x\in A,\\0,&x\notin A,\end{cases}
\qquad \int_{\mathcal S}f(s)\,\mathrm d\delta_x(s)=f(x).
\end{equation*}
$\delta_x$ 是集中于 $x$ 的概率测度，与表示概率单纯形的大写 $\Delta$ 不同。

\textbf{概率测度与经验测度：}$\mathcal B(\mathcal S)$ 表示 Borel 集合族，$\mathcal M(\mathcal S)$ 表示其上的全体测度；概率测度全体记为 $\mathcal M_+^1(\mathcal S):=\{\mu\in\mathcal M(\mathcal S):\mu\ge0,\ \mu(\mathcal S)=1\}$。对 $\bm w\in\Delta^{N-1}$，
\begin{equation*}
\langle\bm w,\delta_{x_{1..N}}\rangle:=\sum_{i=1}^Nw_i\delta_{x_i},\qquad
\hat{\mathbb P}_{N,\bm w;X}:=\sum_{i=1}^Nw_i\delta_{x_i},\qquad
\hat{\mathbb P}_{N;X}:=\frac1N\sum_{i=1}^N\delta_{x_i}.
\end{equation*}
这里 $x_i$ 是样本实现值；权重与测度族的尖括号表示线性组合，其结果仍为测度。

\textbf{样本排列与滤波阶段：}$X_i\in\mathbb R^{d\times1}$ 为样本列向量，
\begin{equation*}
X_{1..N}:=(X_1,\dots,X_N)\in\mathbb R^{d\times N},\qquad
\mathbb P_{X_{1..N}}\text{ 为该样本族的联合分布}.
\end{equation*}
上标 $-$、$+$ 分别表示预测、后验阶段。位置变换中的 $X_k$ 表示按列拼接的粒子矩阵，每列是一个 $n_x$ 维粒子：
\begin{equation*}
\begin{aligned}
&X_k:=x_k^{(1..N)}:=[x_k^{(1)},\dots,x_k^{(N)}]\in\mathbb R^{n_x\times N},\\
&X_k^\pm:=x_k^{\pm,(1..N)}\in\mathbb R^{n_x\times N}.
\end{aligned}
\end{equation*}
$(\bm T_k)_{ij}$ 表示源粒子 $i$ 对目标粒子 $j$ 的贡献系数，故 $X_k^+=X_k^-\bm T_k$，$\bm T_k=N\bm\Pi_k$。概率分布下标中的 $X_k$ 仍表示状态随机变量，由上下文区分。

\textbf{推前、拉回与对偶配对：}对可测映射 $T:\mathcal X\to\mathcal Y$，
\begin{equation*}
\begin{aligned}
&(T_*\mu)(B):=\mu(T^{-1}(B)),\qquad T^*f:=f\circ T,\\
&\langle f,\mu\rangle_{\mathcal S}:=\int_{\mathcal S}f(s)\,\mathrm d\mu(s),\qquad
\langle f,T_*\mu\rangle_{\mathcal Y}=\langle T^*f,\mu\rangle_{\mathcal X}.
\end{aligned}
\end{equation*}
其中 $B\in\mathcal B(\mathcal Y)$，积分对有界可测或相应可积函数使用。$\pi_i$ 表示取第 $i$ 个分量的坐标投影，$\mathbb P_X=X_*\mathbb P$。

\textbf{联合概率矩阵与条件概率矩阵：}在正质量离散支撑上，记
\begin{equation*}
\begin{aligned}
&a_i:=\mathbb P_X(\{x_i\}),\qquad b_j:=\mathbb P_Y(\{y_j\}),\qquad
\Pi_{ij}:=\mathbb P_{(X,Y)}(\{(x_i,y_j)\}),\\
&P_{Y\mid X;i,j}:=\mathbb P_{Y\mid X=x_i}(\{y_j\})=\Pi_{ij}/a_i,\qquad
P_{X\mid Y;j,i}:=\mathbb P_{X\mid Y=y_j}(\{x_i\})=\Pi_{ij}/b_j.
\end{aligned}
\end{equation*}
条件矩阵的行标识条件变量，列标识被条件化变量；$\bm\Pi$ 是联合概率矩阵。$\Pi(\mathbb P_X,\mathbb P_Y)$ 表示具有指定边缘的耦合测度集合，$U(\bm a,\bm b)$ 为其离散矩阵表示，见式~\eqref{cw:eq:coupling-family}、\eqref{cw:eq:transport-polytope}。

\textbf{代价矩阵与矩阵运算：}$C_{ij}:=c(x_i,y_j)$，$\bm C=(C_{ij})$；$\langle\bm A,\bm B\rangle_F:=\sum_{i,j}A_{ij}B_{ij}$ 为实矩阵的 Frobenius 内积。$\operatorname{diag}(\bm u)$ 将向量放在矩阵对角线上，$\odot$、$\oslash$ 表示逐元素乘、除，$\|\bm z\|_1:=\sum_i|z_i|$，$\|\bm z\|$ 表示欧氏范数。

\textbf{乘积与正则化：}$\mu\otimes\nu$ 表示乘积测度，$\mu\ltimes K$ 表示由边缘测度与条件转移核构成的半直积测度，$\otimes_k$ 表示矩阵 Kronecker 积。$D_{\mathrm{KL}}(\mu\|\nu):=\int\log(\mathrm d\mu/\mathrm d\nu)\,\mathrm d\mu$（$\mu\ll\nu$ 时），否则为 $+\infty$；$\varepsilon>0$ 为熵正则化参数，$\tau>0$ 为边缘残差容差。



% 新证明单独维护，保持附录主干简洁。
% !TeX root = ../../../main.tex
\section{无记忆性引理的离散情形证明}\label{cw:app:markov-proof}

\begin{proof}[引理~\ref{cw:lem:markov} 的证明（离散情形）]
沿用引理~\ref{cw:lem:markov} 的模型与独立性条件，设 $X_0$ 和各噪声取有限或可数个值，$p$ 表示相对于计数测度的概率质量函数。
以下仅考虑概率为正的条件取值，求和遍历相应噪声的所有可能取值。

由递推关系，$(X_{0:k-1},Y_{1:k-1})$ 是 $(X_0,N_{x,1:k-1},N_{y,1:k-1})$ 的函数，因此与 $N_{x,k}$ 独立。
固定历史取值后，$X_{k-1}=x_{k-1}$，先按噪声取值分解，再依次使用 $X_k=f(X_{k-1},N_{x,k})$ 与噪声独立性：
\begin{equation*}
\begin{aligned}
&p_{X_k\mid X_{0:k-1},Y_{1:k-1}}(x_k\mid x_{0:k-1},y_{1:k-1})\\
= &\sum_n p_{(X_k,N_{x,k})\mid X_{0:k-1},Y_{1:k-1}}(x_k,n\mid x_{0:k-1},y_{1:k-1})\\
= &\sum_n\mathbb I_{\{f(x_{k-1},n)=x_k\}}p_{N_{x,k}\mid X_{0:k-1},Y_{1:k-1}}(n\mid x_{0:k-1},y_{1:k-1})\\
= &\sum_n\mathbb I_{\{f(x_{k-1},n)=x_k\}}p_{N_{x,k}}(n)\\
= &p_{X_k\mid X_{k-1}}(x_k\mid x_{k-1}).
\end{aligned}
\end{equation*}
最后一步同样利用 $N_{x,k}$ 与 $X_{k-1}$ 独立，故只给定 $X_{k-1}$ 时得到相同的求和式。

同理，$(X_{0:k},Y_{1:k-1})$ 是 $(X_0,N_{x,1:k},N_{y,1:k-1})$ 的函数，因此与 $N_{y,k}$ 独立。由 $Y_k=h(X_k,N_{y,k})$ 得
\begin{equation*}
\begin{aligned}
&p_{Y_k\mid X_{0:k},Y_{1:k-1}}(y_k\mid x_{0:k},y_{1:k-1})\\
= &\sum_n\mathbb I_{\{h(x_k,n)=y_k\}}p_{N_{y,k}}(n)\\
= &p_{Y_k\mid X_k}(y_k\mid x_k).
\end{aligned}
\end{equation*}
这里最后一步利用 $N_{y,k}$ 与 $X_k$ 独立。$k=1$ 时将历史观测视为空，即得引理的两个等式。

一般情形对任意 Borel 集合写条件概率，将上述求和替换为对噪声分布的积分，即得相同的条件分布等式；在正文的密度存在假设下，进一步得到几乎处处成立的密度等式。
\end{proof}


\section{Bayes 更新公式的证明}\label{cw:app:proofs}

\begin{theorem}[Bayes 更新公式]\label{cw:thm:bayes}

在本文约定与模型条件下，对 $k\geq1$，给定历史观测 $y_{1:k-1}$ 并收到当前观测 $y_k$ 后，状态后验密度满足
\begin{equation}\label{cw:eq:15}
\begin{aligned}
&p_{X_k\mid Y_{1:k}}(x_k\mid y_{1:k})\\
= &\frac{p_{Y_k\mid X_k}(y_k\mid x_k)p_{X_k\mid Y_{1:k-1}}(x_k\mid y_{1:k-1})}
{\displaystyle\int_{\mathcal{X}}p_{Y_k\mid X_k}(y_k\mid x)p_{X_k\mid Y_{1:k-1}}(x\mid y_{1:k-1})\,\mathrm{d}v_X(x)}.
\end{aligned}
\end{equation}

\end{theorem}

\begin{proof}

先固定历史观测 $y_{1:k-1}$。收到 $y_k$ 后，要求的是在这些历史观测的基础上，再给定 $Y_k=y_k$ 时的状态密度。按条件密度公式，将条件联合密度除以新观测的条件边缘密度：
\begin{equation}\label{cw:eq:16}
\begin{aligned}
&p_{X_k\mid Y_{1:k}}(x_k\mid y_{1:k})\\
= &\frac{p_{(X_k,Y_k)\mid Y_{1:k-1}}(x_k,y_k\mid y_{1:k-1})}
{p_{Y_k\mid Y_{1:k-1}}(y_k\mid y_{1:k-1})}.
\end{aligned}
\end{equation}

对分子使用密度的乘法公式，按“先取状态，再给定状态取观测”的顺序分解。两个因子都保留历史观测条件，其中状态因子正是预测步骤已求得的密度：
\begin{equation}\label{cw:eq:17}
\begin{aligned}
&p_{(X_k,Y_k)\mid Y_{1:k-1}}(x_k,y_k\mid y_{1:k-1})\\
= &p_{Y_k\mid X_k,Y_{1:k-1}}(y_k\mid x_k,y_{1:k-1})
p_{X_k\mid Y_{1:k-1}}(x_k\mid y_{1:k-1}).
\end{aligned}
\end{equation}
将这一分解代入上一式，就是以历史观测为条件的 Bayes 公式。

接着应用无记忆性引理中的观测条件独立性。给定当前状态 $X_k$ 后，历史观测不再影响 $Y_k$ 的条件分布，因此观测因子可以直接使用模型给出的似然：
\begin{equation}\label{cw:eq:18}
p_{Y_k\mid X_k,Y_{1:k-1}}(y_k\mid x_k,y_{1:k-1})
=p_{Y_k\mid X_k}(y_k\mid x_k).
\end{equation}
这里仅消去了观测因子中的历史条件，预测密度仍依赖历史观测。

分母由条件联合密度对状态变量积分得到，这是条件全概率公式。再代入上面的乘法分解与条件独立关系，就将观测的预测密度写成似然关于状态预测分布的平均：
\begin{equation}\label{cw:eq:19}
\begin{aligned}
&p_{Y_k\mid Y_{1:k-1}}(y_k\mid y_{1:k-1})\\
= &\int_{\mathcal{X}}p_{(X_k,Y_k)\mid Y_{1:k-1}}(x,y_k\mid y_{1:k-1})\,\mathrm{d}v_X(x)\\
= &\int_{\mathcal{X}}p_{Y_k\mid X_k}(y_k\mid x)
p_{X_k\mid Y_{1:k-1}}(x\mid y_{1:k-1})\,\mathrm{d}v_X(x).
\end{aligned}
\end{equation}
其中 $x$ 是积分变量，覆盖整个状态空间。按前述条件，这个分母有限且严格为正，可以用于归一化。

最后，将得到的分子与分母代回最初的条件密度公式，即得定理中的更新式。分母是似然加权后预测密度的总积分，因此相除后得到积分为 $1$ 的后验密度。
\end{proof}


\section{离散 Monge 映射的存在条件}\label{cw:app:monge}

\textbf{原子质量守恒：}设 $x_1,\dots,x_m$ 各自互异，$y_1,\dots,y_n$ 各自互异，$a_i,b_j>0$。若 $T_*\mathbb P_X=\mathbb P_Y$，则
\begin{equation}\label{cw:eq:monge-atom-mass}
\begin{aligned}
b_j&=\mathbb E_{Y\sim\mathbb P_Y}[\mathbb I_{\{y_j\}}(Y)]
=\mathbb E_{X\sim\mathbb P_X}[\mathbb I_{\{y_j\}}(T(X))]\\
&=\int_{\mathcal X}\mathbb I_{\{y_j\}}(T(x))\,
\mathrm d\!\left(\sum_{i=1}^m a_i\delta_{x_i}\right)(x)
=\sum_{\{i:\,T(x_i)=y_j\}}a_i.
\end{aligned}
\end{equation}
每个正质量源点必然映入目标支撑，因为
\begin{equation}\label{cw:eq:monge-support}
\mathbb P_Y(\{T(x_i)\})
=\mathbb P_X(T^{-1}(\{T(x_i)\}))\ge a_i>0.
\end{equation}
因此其条件概率矩阵为逐行单点选择，联合概率矩阵则带有源权重：
\begin{equation}\label{cw:eq:monge-conditional}
\begin{aligned}
&P_{Y\mid X;i,j}=\mathbb I_{\{y_j\}}(T(x_i))\in\{0,1\},\qquad
\sum_{j=1}^n P_{Y\mid X;i,j}
=\delta_{T(x_i)}\!\left(\bigcup_{j=1}^n\{y_j\}\right)=1,\\
&\bm P_{Y\mid X}\mathbf1_n=\mathbf1_m,\qquad
\bm P_{Y\mid X}^{\mathrm T}\bm a=\bm b,\qquad
\bm\Pi=\operatorname{diag}(\bm a)\bm P_{Y\mid X}.
\end{aligned}
\end{equation}
同一质量关系也可以完全用推前与拉回对偶表达：
\begin{equation}\label{cw:eq:monge-duality}
\begin{aligned}
b_j&=\langle\mathbb I_{\{y_j\}},\mathbb P_Y\rangle_{\mathcal Y}
=\langle\mathbb I_{\{y_j\}},T_*\mathbb P_X\rangle_{\mathcal Y}
=\langle T^*\mathbb I_{\{y_j\}},\mathbb P_X\rangle_{\mathcal X}\\
&=\left\langle T^*\mathbb I_{\{y_j\}},
\langle\bm a,\delta_{x_{1..m}}\rangle\right\rangle_{\mathcal X}
=\sum_{i=1}^m a_iP_{Y\mid X;i,j}.
\end{aligned}
\end{equation}

% 保留定理陈述的最小空间，证明可自然跨页，不预留半页。
\Needspace{5\baselineskip}
\begin{theorem}[均匀离散分布的整除条件]\label{cw:thm:monge-divisibility}
设源、目标分布分别均匀分布于 $m$ 个与 $n$ 个互异支撑点。存在可测映射 $T$ 满足 $T_*\mathbb P_X=\mathbb P_Y$，当且仅当 $n\mid m$。
\end{theorem}
\begin{proof}
必要性：由式~\eqref{cw:eq:monge-atom-mass}，每个目标点的原像必须含有 $m/n$ 个源点：
\begin{equation}\label{cw:eq:monge-divisibility}
\frac1n=\frac1m\#\{i:T(x_i)=y_j\}
\quad\Longrightarrow\quad\frac mn\in\mathbb N.
\end{equation}
充分性：若 $m=nq$，将每 $q$ 个源点分配给一个目标点。用 Kronecker 积写为
\begin{equation}\label{cw:eq:monge-construction}
\begin{aligned}
&\bm P_{Y\mid X}=\bm I_n\otimes_k\mathbf1_q\in\{0,1\}^{m\times n},
\qquad\bm\Pi=\frac1m\bm P_{Y\mid X},\\
&\bm P_{Y\mid X}\mathbf1_n
=(\bm I_n\otimes_k\mathbf1_q)(\mathbf1_n\otimes_k1)
=\mathbf1_n\otimes_k\mathbf1_q=\mathbf1_m,\\
&\bm P_{Y\mid X}^{\mathrm T}\frac{\mathbf1_m}{m}
=\frac{(\bm I_n\otimes_k\mathbf1_q^{\mathrm T})(\mathbf1_n\otimes_k\mathbf1_q)}m
=\frac q m\mathbf1_n=\frac1n\mathbf1_n.
\end{aligned}
\end{equation}
以矩阵每行的唯一非零元素确定 $T(x_i)$，并在有限源支撑之外任取常值延拓，即得到所需可测映射。
\end{proof}
